% !TEX root = model_writeup_quantitative_nospace.tex
\documentclass[11pt]{article}
\usepackage[fleqn]{amsmath}
\usepackage{amssymb}
\usepackage{amsthm}
\usepackage{geometry}
\usepackage{enumitem}
\usepackage{booktabs}
\usepackage[authoryear]{natbib}
\usepackage[colorlinks=true,citecolor=blue,linkcolor=blue,urlcolor=blue]{hyperref}

\geometry{margin=1in}
\setlength{\parskip}{0.6em}
\setlength{\parindent}{0pt}

\newtheorem{definition}{Definition}
\newtheorem{remark}{Remark}

\newcommand{\R}{\mathbb{R}}
\newcommand{\ind}{\mathbb{I}}
\newcommand{\E}{\mathbb{E}}
\newcommand{\Nset}{\mathcal N}
\newcommand{\Dset}{\mathcal D}
\newcommand{\Zset}{\mathcal Z}
\newcommand{\Sset}{\mathcal S}
\newcommand{\Oset}{\mathcal O}
\newcommand{\Tg}{\mathcal T^g}
\newcommand{\Bcal}{\mathcal B}

\title{Fertility and Housing Tenure Choice:\\Quantitative One-Market Model Writeup}
\author{Tommaso De Santo}
\date{June 2026}

\begin{document}
\maketitle

\section{Model}

This version rewrites the May 2026 model as a non-spatial quantitative
lifecycle economy. The May draft had residential locations, local amenities,
bilateral moving wedges, and location-specific housing supply. The current
model removes that spatial layer. There is one housing-services market, one
purchase price, and one rental user cost. What remains is the channel that is
central for the quantitative exercise: children require goods and housing
space, family-suitable housing is easier to access through ownership than
through renting, and the ownership margin is shaped by down-payment constraints,
transaction costs, bequest motives, and tax incentives to retain housing.

The object is a stationary overlapping-generations economy. Adult population
mass is normalized in the benchmark, so fertility changes the composition of
households and their housing demand rather than mechanically changing aggregate
city scale. This normalization keeps the quantitative model focused on the
family-housing margin. A separate population-scale closure can be added later,
but it is not needed to define the current equilibrium.

The model combines endogenous fertility in a lifecycle problem
\citep{sommerFertilityChoiceLife2016,doepkeBargainingBabiesTheory2019} with
housing tenure choice, collateral constraints, and owner-occupied housing
\citep{hendersonModelHousingTenure1983,sommerEquilibriumEffectFundamentals2013,sommerImplicationsUSTax2018}.
Relative to the May draft, the Rosen--Roback spatial layer is removed: there
are no local amenities, no bilateral moving wedges, and no location-specific
housing markets in the baseline quantitative economy.

\subsection{Environment}

\textbf{Agents.}
Households live for ages $a=1,\ldots,J$. They work until age $J_R$ and are
retired afterward. At the beginning of a period, the core individual state is
\begin{equation}
    x=(b,z,d,a,n,s).
    \label{eq:state}
\end{equation}
Here $b$ is liquid wealth, $z\in\Zset$ is idiosyncratic labor productivity,
$d\in\Dset$ is tenure and housing position, $a$ is age, $n$ is completed
fertility, and $s\in\Sset$ is the child-at-home state. The tenure set is
\begin{equation*}
    \Dset=\{R,O_1,\ldots,O_K\}.
\end{equation*}
$R$ denotes renting. $O_k$ denotes ownership of a house with physical size
$H_k$, where $H_1<\cdots<H_K$. Renters choose rental space continuously up to a
cap $\bar h^R$, while owners choose from the discrete ladder $\{H_k\}_{k=1}^K$.
The ladder is the reduced-form way the quantitative model captures the fact
that large family-suitable units are disproportionately owner-occupied.

When the capital-gains module is active, the owner state is augmented by a tax
basis $B^p$. The baseline implementation can be recovered by setting the
capital-gains tax rate to zero, in which case $B^p$ drops out of the household
problem.

\textbf{Income.}
Working-age households receive after-tax labor income
\begin{equation*}
    y_a(z)=(1-\tau^{\mathrm{pay}})w e_a z,
    \qquad a\leq J_R,
\end{equation*}
where $e_a$ is the deterministic lifecycle efficiency profile and $z$ follows
the Markov transition matrix $\Pi^z(z'|z)$. Retired households receive pension
income $p^{\mathrm{ret}}_a$. The quantitative model treats the income process
as a calibrated household primitive rather than as the source of a spatial
wage gradient.

\textbf{Preferences.}
Let $m(n,s)$ be the number of dependent children currently living at home.
Completed fertility $n$ and current dependents $m(n,s)$ differ because children
eventually leave the household. Child-related goods and housing needs are
\begin{equation}
\begin{aligned}
\bar c(n,s)
    &= \bar c_0+\bar c_n m(n,s),\\
\bar h(n,s)
    &= \bar h_0
       +\ind\{m(n,s)\geq 1\}\bar h_{\mathrm{jump}}
       +\bar h_n m(n,s).
\end{aligned}
\label{eq:child_needs}
\end{equation}
The first child can therefore create a discrete housing-space requirement, and
additional children raise required space linearly.

The flow utility aggregator is Stone--Geary over nondurable consumption and
housing services:
\begin{equation}
u(c,\widetilde h;n,s)
=
\frac{
\left[
\left(c-\bar c(n,s)\right)^\alpha
\widetilde h^{\,1-\alpha}
\right]^{1-\sigma}
}{1-\sigma}
\psi_{\mathrm{child}}m(n,s),
\label{eq:flow_utility}
\end{equation}
with utility equal to $-\infty$ if consumption or housing services fall below
the required levels. The object $\widetilde h$ is residual housing service
after child space needs. For renters and owners,
\begin{equation}
\widetilde h =
\begin{cases}
h^R-\bar h(n,s), & d=R,\\
\chi\left[H_k-\bar h(n,s)\right], & d=O_k.
\end{cases}
\label{eq:effective_housing_services}
\end{equation}
The parameter $\chi$ is an owner-occupancy service premium applied to residual
owner housing services. In market clearing, owner demand is the physical
quantity $H_k$, not the service quantity $\chi H_k$.

The key fertility mechanism is already visible in \eqref{eq:flow_utility} and
\eqref{eq:effective_housing_services}. Children do not only require goods.
They also absorb scarce housing space. If a household is constrained by the
rental cap, the ownership menu, or the down-payment requirement, the relevant
price of a child is the household's shadow value of space, not only the market
rent.

\textbf{Bequests.}
At death, households receive warm-glow bequest utility. Let $W^B$ denote
bequeathable wealth after any liquidation costs and taxes. Bequest utility is
\begin{equation}
\Bcal(W^B,n)
=
\theta_0(1+\theta_n n)
\frac{
\left(\theta_1+\max\{W^B,0\}\right)^{1-\sigma}
-\theta_1^{1-\sigma}
}{1-\sigma}.
\label{eq:bequest}
\end{equation}
The parameter $\theta_0$ governs the overall bequest motive, $\theta_n$ lets
completed fertility scale the bequest motive, and $\theta_1$ controls curvature
near zero wealth.

\textbf{Housing market.}
There is one aggregate housing-services market. The rental user cost is $q$ and
the asset price of housing is $P$. They satisfy
\begin{equation}
    q=(r+\delta+\tau^p)P,
    \label{eq:user_cost}
\end{equation}
where $r$ is the risk-free interest rate, $\delta$ is maintenance/depreciation,
and $\tau^p$ is the recurrent property-tax rate. Aggregate housing supply is
upward sloping:
\begin{equation}
    H^S(q)=H_0\left(\frac{q}{\bar q}\right)^\eta.
    \label{eq:housing_supply}
\end{equation}
The scale $H_0$ fixes the benchmark level of housing services and $\eta$
governs the construction response.

\textbf{Financing and tenure institutions.}
Liquid wealth earns gross return $R=1+r$. Renters cannot borrow. A household
that originates an owner position $O_k$ can finance at most share $\phi$ of the
house value. Equivalently, the required down payment is
\begin{equation}
    (1-\phi)PH_k.
    \label{eq:down_payment}
\end{equation}
The parameter $\phi$ is the financed share, so a larger $\phi$ relaxes the
ownership constraint.

Selling a house is costly. If a household sells a house of size $H_k$, sale
proceeds before capital-gains taxes are
\begin{equation*}
    (1-\psi^s)PH_k,
\end{equation*}
where $\psi^s$ is the proportional sale or transaction wedge. A household that
keeps its current house is not forced to satisfy a new origination constraint
after a price change. This distinction matters for old owners: the model can
generate retention because selling is costly, because current owners already
hold large units, and because tax rules can make holding more attractive than
realizing gains.

\textbf{Capital gains and step-up basis.}
The capital-gains block is modular. Let $B^p$ be the tax basis attached to an
owner-occupied house and let $\bar g$ be the exempt gain. If the house is sold
during life, the tax liability is
\begin{equation}
    \Tg(PH_k,B^p)
    =
    \tau^g\max\{PH_k-B^p-\bar g,0\}.
    \label{eq:capital_gains_tax}
\end{equation}
Sale proceeds net of transaction costs and capital-gains taxes are therefore
\begin{equation}
    \Omega(H_k,B^p)
    =
    (1-\psi^s)PH_k-\Tg(PH_k,B^p).
    \label{eq:sale_proceeds}
\end{equation}
Under step-up basis, accrued gains are not taxed at death. The tax basis for
heirs is reset to the market value of the inherited house. This creates a
retention incentive for older owners with large unrealized gains: selling
converts the gain into a current tax liability, while holding until death can
forgive that liability. Setting $\tau^g=0$ or $B^p=PH_k$ collapses this block to
the current no-capital-gains baseline.

\subsection{Household Problem}

Within a period, the household first faces the fertility choice if eligible,
then chooses tenure and housing position, and finally chooses consumption,
rental space, and saving. The order is useful because fertility changes
housing needs before the household chooses whether to rent, keep an existing
home, or buy a different owner rung.

\textbf{Housing adjustment.}
Let $H(R)=0$ and $H(O_k)=H_k$. Given current tenure $d$ and next tenure $d'$,
post-adjustment liquid wealth is
\begin{equation}
\widetilde b(b,d,d')
=
b
+\ind\{d\in\Oset,\ d'\neq d\}\Omega(H(d),B^p)
-\ind\{d'\in\Oset,\ d'\neq d\}PH(d'),
\label{eq:adjusted_wealth}
\end{equation}
where $\Oset=\{O_1,\ldots,O_K\}$. If the capital-gains block is inactive, the
argument $B^p$ is suppressed and $\Omega(H(d),B^p)=(1-\psi^s)PH(d)$.

The feasible next-tenure set is
\begin{equation}
\Dset(b,d)
=
\left\{
d'\in\Dset:
\widetilde b(b,d,d')\geq \underline b(d'|b,d)
\right\}.
\label{eq:feasible_tenure}
\end{equation}
For renters and new owner positions,
\begin{equation*}
\underline b(R|b,d)=0,
\qquad
\underline b(O_k|b,d)=-\phi PH_k.
\end{equation*}
For an incumbent owner who keeps the same house, the lower bound is
\begin{equation}
    \underline b(O_k|b,O_k)=\min\{b,-\phi PH_k\}.
    \label{eq:incumbent_ltv}
\end{equation}
Thus existing owners are not forced to delever mechanically when prices move,
but households that buy or upsize must satisfy the origination down-payment
constraint.

\textbf{Branch values.}
Let $V_{a+1}$ be next period's value. The continuation value after income risk
and child aging is
\begin{equation}
\bar V_{a+1}(b',d',z,n,s)
=
\sum_{z'\in\Zset}\sum_{s'\in\Sset}
\Pi^z(z'|z)\Pi^c(s'|s,n)
V_{a+1}(b',z',d',a+1,n,s').
\label{eq:continuation}
\end{equation}
At the terminal date, $V_{J+1}(b,z,d,J+1,n,s)=\Bcal(W^B,n)$.

Conditional on choosing to rent, the household solves
\begin{equation}
\begin{aligned}
V^R_a(\widetilde b,z,n,s)
=
\max_{c,h^R,b'}\quad
&u(c,h^R-\bar h(n,s);n,s)
+\beta \bar V_{a+1}(b',R,z,n,s)\\
\text{s.t.}\quad
&c+q h^R+b'=R\widetilde b+y_a(z),\\
&0\leq h^R\leq \bar h^R,\qquad b'\geq 0.
\end{aligned}
\label{eq:renter_problem}
\end{equation}
The rental cap is a primitive family-housing access margin: when it binds, the
household's marginal value of space can exceed the market rent.

Conditional on owning rung $O_k$, the household solves
\begin{equation}
\begin{aligned}
V^{O_k}_a(\widetilde b,z,n,s)
=
\max_{c,b'}\quad
&u\!\left(c,\chi[H_k-\bar h(n,s)];n,s\right)
+\beta \bar V_{a+1}(b',O_k,z,n,s)\\
\text{s.t.}\quad
&c+(\delta+\tau^p)PH_k+b'=R\widetilde b+y_a(z),\\
&b'\geq \underline b^+ (O_k).
\end{aligned}
\label{eq:owner_problem}
\end{equation}
The lower bound $\underline b^+(O_k)$ is the continuation borrowing limit
associated with the owner position. The owner pays recurrent property taxes and
maintenance through the service cost $(\delta+\tau^p)PH_k$.

\textbf{Tenure and house-size choice.}
For each feasible next tenure $d'$, define the deterministic branch value
\begin{equation}
\widetilde V^T_a(d'|b,z,d,n,s)
=
\begin{cases}
V^R_a(\widetilde b(b,d,R),z,n,s), & d'=R,\\
V^{O_k}_a(\widetilde b(b,d,O_k),z,n,s), & d'=O_k.
\end{cases}
\label{eq:tenure_branch_value}
\end{equation}
Tenure and owner-rung choices are smoothed with Type-I extreme-value shocks of
scale $\kappa_T$. The probability of choosing $d'$ is
\begin{equation}
\pi^T_a(d'|b,z,d,n,s)
=
\frac{
\exp\left(\widetilde V^T_a(d'|b,z,d,n,s)/\kappa_T\right)
}{
\sum_{\hat d\in\Dset(b,d)}
\exp\left(\widetilde V^T_a(\hat d|b,z,d,n,s)/\kappa_T\right)
},
\label{eq:tenure_probs}
\end{equation}
and the inclusive tenure value is
\begin{equation}
V^T_a(b,z,d,n,s)
=
\kappa_T
\log
\sum_{\hat d\in\Dset(b,d)}
\exp\left(\widetilde V^T_a(\hat d|b,z,d,n,s)/\kappa_T\right).
\label{eq:tenure_iv}
\end{equation}
As $\kappa_T$ approaches zero, tenure choice converges to the deterministic
maximum. With positive $\kappa_T$, the model allows small residual dispersion in
tenure and size choices.

\textbf{Fertility.}
During the fertile window
\begin{equation*}
    a\in\{A_f^{\mathrm{start}},\ldots,A_f^{\mathrm{end}}\},
\end{equation*}
a childless household chooses $n'\in\Nset=\{0,1,\ldots,\bar n\}$. The option
$n'=0$ means no birth in the current period. A positive option fixes completed
fertility forever and moves the household into the first dependent-child state.
Let
\begin{equation*}
    s(n')=
    \begin{cases}
    0, & n'=0,\\
    1, & n'>0.
    \end{cases}
\end{equation*}
The value of fertility option $n'$ is
\begin{equation}
\widetilde V^F_a(n'|b,z,d)
=
V^T_a(b,z,d,n',s(n')).
\label{eq:fert_branch_value}
\end{equation}
Fertility choices are smoothed with Type-I extreme-value shocks of scale
$\kappa_F$:
\begin{equation}
\pi^F_a(n'|b,z,d)
=
\frac{
\exp\left(\widetilde V^F_a(n'|b,z,d)/\kappa_F\right)
}{
\sum_{\hat n\in\Nset}
\exp\left(\widetilde V^F_a(\hat n|b,z,d)/\kappa_F\right)
}.
\label{eq:fert_probs}
\end{equation}
The household value at a childless fertile state is
\begin{equation}
V_a(b,z,d,a,0,0)
=
\kappa_F
\log
\sum_{\hat n\in\Nset}
\exp\left(\widetilde V^F_a(\hat n|b,z,d)/\kappa_F\right).
\label{eq:fert_iv}
\end{equation}
Outside the fertile window, and after a positive fertility choice has already
been made, fertility is degenerate and $V_a(b,z,d,a,n,s)=V^T_a(b,z,d,n,s)$.

This block should be interpreted as a completed-family choice in the current
quantitative model. It is not a sequential second-birth hazard model. The
sequential parity version is a natural extension, but it requires a larger
fertility state space and different empirical moments.

\subsection{Stationary Equilibrium}

The equilibrium is stationary in prices, household policies, and the normalized
distribution of adult households. The aggregate adult mass is normalized to one.
This normalization means the model clears the housing market for a stationary
composition, not for an endogenous city size.

\begin{definition}[Stationary equilibrium]
A stationary equilibrium is a collection
\begin{equation*}
\begin{aligned}
\mathcal X^\ast=\big(&P,q,
\{V_a,V^R_a,V^{O_k}_a,V^T_a,\widetilde V^T_a,\widetilde V^F_a\}_{a,k},
\{c,h^R,b',d',n'\},\\
&\{\pi^T,\pi^F\},
g,H^D,H^S
\big)
\end{aligned}
\end{equation*}
such that:
\begin{enumerate}[leftmargin=*]
    \item \textbf{Household optimization.} Given prices, taxes, and the income
    process, value functions solve the household problem in
    \eqref{eq:renter_problem}--\eqref{eq:owner_problem}, with tenure probabilities
    \eqref{eq:tenure_probs}, fertility probabilities \eqref{eq:fert_probs}, and
    terminal bequest utility \eqref{eq:bequest}. Policies for consumption, rental
    space, saving, tenure, owner rung, and fertility are associated optimal
    policies.

    \item \textbf{Stationary distribution.} The normalized adult distribution
    $g$ is invariant under the transition induced by optimal policies, tenure
    and fertility probabilities, the income transition, child aging, deterministic
    aging, and death. Adult mass is normalized:
    \begin{equation*}
        \int dg(x)=1.
    \end{equation*}
    New adult households enter at the initial age with the benchmark distribution
    over income, assets, tenure, fertility, and child state.

    \item \textbf{Housing demand.} Aggregate physical housing demand is
    \begin{equation}
    H^D
    =
    \int
    \left[
    \ind\{d=R\}h^R(x)
    +
    \sum_{k=1}^K\ind\{d=O_k\}H_k
    \right]dg(x).
    \label{eq:housing_demand}
    \end{equation}
    Housing demand counts physical space. It does not multiply owner-occupied
    units by the service premium $\chi$.

    \item \textbf{Housing market clearing and user cost.} The aggregate housing
    market clears:
    \begin{equation}
        H^D=H^S(q),
        \qquad
        q=(r+\delta+\tau^p)P.
        \label{eq:market_clearing}
    \end{equation}

    \item \textbf{Government accounting.} Payroll taxes finance the pension
    profile if the pension block is closed inside the model. Property taxes and
    capital-gains taxes are rebated or absorbed by the government according to
    the counterfactual experiment. When the capital-gains module is inactive,
    $\tau^g=0$ and the tax basis state is irrelevant.
\end{enumerate}
\end{definition}

\section{Quantitative Content}

The quantitative model is designed to connect housing access to fertility
through a small number of observable margins. The calibration should discipline
the object that appears in the mechanism: the shadow price of family-suitable
housing faced by young households.

\subsection{Moment Blocks}

The paper-facing calibration table should be organized around parameter blocks,
not around the order in which the code computes moments:
\begin{enumerate}[leftmargin=*]
    \item \textbf{Fertility.} Completed fertility, total fertility, and
    childlessness discipline the net value of children, the curvature of the
    fertility choice, and the child cost parameters.

    \item \textbf{Lifecycle tenure.} Ownership by age and old-age ownership
    discipline down-payment constraints, transaction costs, and the bequest or
    retention motives.

    \item \textbf{Family housing demand.} Housing-size changes around children,
    rooms or housing services by child status, and the renter-owner size gap
    discipline $\bar h_{\mathrm{jump}}$, $\bar h_n$, the renter cap $\bar h^R$,
    and the owner size ladder.

    \item \textbf{Market clearing and supply.} The level of housing services and
    the user cost identify the supply scale $H_0$ and the price response
    governed by $\eta$.

    \item \textbf{Wealth and bequests.} Assets by age, housing wealth at old
    ages, and bequest-related moments discipline the warm-glow bequest
    parameters and the strength of retention at the end of life.
\end{enumerate}

This organization keeps the calibration tied to the economics. Fertility
moments alone do not identify the housing block. Housing-size moments alone do
not identify the value of children. Old-age ownership is informative about
retention only once the transaction, tax, and bequest channels are stated
explicitly.

\subsection{Counterfactuals}

The natural counterfactuals change one primitive access margin at a time:
\begin{enumerate}[leftmargin=*]
    \item relax the down-payment constraint by increasing $\phi$;
    \item expand rental access to family-sized housing by increasing $\bar h^R$;
    \item reduce transaction or sale wedges $\psi^s$;
    \item remove or change capital-gains taxation and step-up basis;
    \item increase housing supply elasticity $\eta$ or the supply scale $H_0$;
    \item change child-related housing needs $\bar h_{\mathrm{jump}}$ and
    $\bar h_n$ to measure the importance of the space requirement itself.
\end{enumerate}

These experiments should be interpreted as general-equilibrium exercises in the
one housing market. They change prices through \eqref{eq:market_clearing}, but
they do not include migration, spatial sorting, or city-scale feedbacks.

\section{Relation to the May Spatial Draft}

The May model had two locations, local amenities, local housing supply, and an
outside-option entry system for city size. Those objects are deliberately absent
from the current quantitative model. The mechanism is now national or
one-market rather than spatial: the relevant scarcity is access to family-sized
housing services, especially through the tenure system, not the allocation of
families across center and periphery.

The substantive elements retained from the May draft are the lifecycle problem,
the renter-owner distinction, the owner size ladder, the rental-size cap, the
down-payment constraint, child housing needs, bequest motives, and housing
market clearing. The main new element relative to the May presentation is the
capital-gains and step-up-basis module. That module is useful because it creates
a direct old-owner retention wedge: an older household with large unrealized
gains may hold a large house even when its current housing need is low, because
selling triggers taxes that can be avoided by holding until death.

The current draft should therefore not claim that old occupancy is inefficient
by itself. A large house occupied by an old owner is inefficient only relative
to a stated friction or policy wedge: transaction costs, capital-gains taxes,
step-up basis, property-tax preferences, or a missing market that prevents young
families from accessing equivalent space. The quantitative task is to measure
how much those wedges raise young households' shadow value of family housing
and how much that shadow value depresses fertility.

\appendix
\section{Optional Population Closure}

The benchmark normalization fixes adult mass at one. If the paper later needs a
scale channel, the least disruptive extension is a renewal equation for young
adult entry. Let $E_0$ be the mass of new adult entrants required to reproduce
one unit of the stationary age distribution, and let $B_0$ be the mature-child
flow produced by one unit of adult mass. With outside-born potential entrant
flow $M$ and retention probability $\rho$, stationary scale would satisfy
\begin{equation}
    S E_0 = M+\rho S B_0,
    \qquad
    S=\frac{M}{E_0-\rho B_0}.
    \label{eq:renewal_scale}
\end{equation}
This adds a scale response to fertility without reintroducing locations. It is
not part of the baseline quantitative model because the current calibration
targets the composition of households and housing choices, not city population
growth.

\bibliographystyle{apalike}
\bibliography{main_bib,lit_review_extra}

\end{document}
